\documentclass{article}

\usepackage{fancyhdr}
\usepackage{extramarks}
\usepackage{amsmath}
\usepackage{amsthm}
\usepackage{amsfonts}
\usepackage{enumitem}
\usepackage{tikz}
\usepackage[plain]{algorithm}
\usepackage{algpseudocode}

\newcommand{\I}{\mathcal{I}}
\newcommand{\F}{\mathbb{F}}
\newcommand{\E}{\mathbb{E}}
\newcommand{\R}{\mathbb{R}}
\renewcommand{\P}{\mathbb{P}}
\usetikzlibrary{automata,positioning}

%
% Basic Document Settings
%

\topmargin=-0.45in
\evensidemargin=0in
\oddsidemargin=0in
\textwidth=6.5in
\textheight=9.0in
\headsep=0.25in

\linespread{1.1}

\pagestyle{fancy}
\lhead{\hmwkAuthorName}
\chead{\hmwkClass\ (\hmwkClassInstructor\ \hmwkClassTime): \hmwkTitle}
\rhead{\firstxmark}
\lfoot{\lastxmark}
\cfoot{\thepage}

\renewcommand\headrulewidth{0.4pt}
\renewcommand\footrulewidth{0.4pt}

\setlength\parindent{0pt}


\newtheorem{theorem}{Theorem}
% 
% Create Problem Sections
%


\newcommand{\enterProblemHeader}[1]{
    \nobreak\extramarks{}{}\nobreak{}
    \nobreak\extramarks{}\nobreak{}
}

\newcommand{\exitProblemHeader}[1]{
    \nobreak\extramarks{}\nobreak{}
    \stepcounter{#1}
    \nobreak\extramarks{}\nobreak{}
  }

  
\newcommand{\enterExerciseHeader}[1]{
    \nobreak\extramarks{}{}\nobreak{}
    \nobreak\extramarks{}\nobreak{}
}

\newcommand{\exitExerciseHeader}[1]{
    \nobreak\extramarks{}\nobreak{}
    \stepcounter{#1}
    \nobreak\extramarks{}\nobreak{}
}

\setcounter{secnumdepth}{0}
\newcounter{partCounter}
\newcounter{homeworkProblemCounter}
\setcounter{homeworkProblemCounter}{1}

\newcounter{homeworkExerciseCounter}
\setcounter{homeworkExerciseCounter}{1}
\nobreak\extramarks{}\nobreak{}
\nobreak\extramarks{}\nobreak{}

%
% Homework Problem Environment
%
% This environment takes an optional argument. When given, it will adjust the
% problem counter. This is useful for when the problems given for your
% assignment aren't sequential. See the last 3 problems of this template for an
% example.
%
\newenvironment{homeworkProblem}[1][-1]{
    \ifnum#1>0
        \setcounter{homeworkProblemCounter}{#1}
    \fi
    \section{Problem \arabic{homeworkProblemCounter}}
    \setcounter{partCounter}{1}
    \enterProblemHeader{homeworkProblemCounter}
}{
    \exitProblemHeader{homeworkProblemCounter}
}

  \newenvironment{homeworkExercise}[1][-1]{
    \ifnum#1>0
        \setcounter{homeworkExerciseCounter}{#1}
    \fi
    \section{Exercise \arabic{homeworkExerciseCounter}}
    \setcounter{partCounter}{1}
    \enterExerciseHeader{homeworkExerciseCounter}
}{
    \exitExerciseHeader{homeworkExerciseCounter}
}

%
% Homework Details
%   - Title
%   - Due date
%   - Class
%   - Section/Time
%   - Instructor
%   - Author
%

\newcommand{\hmwkTitle}{Homework\ \#2}
\newcommand{\hmwkAssignedDate}{February 12, 2017}
\newcommand{\hmwkDueDate}{February 26, 2017, before class}
\newcommand{\hmwkClass}{Design and Analysis of Algorithms, CS 6550}
\newcommand{\hmwkClassTime}{}%Section A}
\newcommand{\hmwkClassInstructor}{Professor Jamie Morgenstern}
\newcommand{\hmwkAuthorName}{}%\textbf{Josh Davis} \and \textbf{Davis Josh}}

%
% Title Page
%

\title{
    \vspace{2in}
    \textmd{\textbf{\hmwkClass:\ \hmwkTitle}}\\
    \normalsize\vspace{0.1in}\small{Due\ on\ \hmwkDueDate\ at 3:10pm}\\
    \vspace{0.1in}\large{\textit{\hmwkClassInstructor\ \hmwkClassTime}}
    \vspace{3in}
}

\author{\hmwkAuthorName}
\date{}

\renewcommand{\part}[1]{\textbf{\large Part \Alph{partCounter}}\stepcounter{partCounter}\\}

%
% Various Helper Commands
%

% Useful for algorithms
\newcommand{\alg}[1]{\textsc{\bfseries \footnotesize #1}}

% For derivatives
\newcommand{\deriv}[1]{\frac{\mathrm{d}}{\mathrm{d}x} (#1)}

% For partial derivatives
\newcommand{\pderiv}[2]{\frac{\partial}{\partial #1} (#2)}

% Integral dx
\newcommand{\dx}{\mathrm{d}x}

% Alias for the Solution section header
\newcommand{\solution}{\textbf{\large Solution}}

% Probability commands: Expectation, Variance, Covariance, Bias
\newcommand{\Var}{\mathrm{Var}}
\newcommand{\Cov}{\mathrm{Cov}}
\newcommand{\Bias}{\mathrm{Bias}}

\begin{document}

% \maketitle

\pagebreak

{\bf Homework Out: \hmwkAssignedDate}\\
{\bf Due Date:  \hmwkDueDate}\\

The HW contains some exercises (fairly simple problems to check you
are on board with the concepts; don’t submit your solutions), and
problems (for which you should submit your solutions, and which will
be graded). Some problems have sub-parts that are exercises.  For this
problem set, it’s OK to work with others. (Groups of 2, maybe 3 max.)
That being said, please think about the problems yourself before
talking to others. Please cite all sources you use, and people you
work with. The expectation is that you try and solve these problems
yourself, rather than looking online explicitly for
answers. Submissions due at beginning of class on the due date.
Please check the Piazza for details on submitting your \emph{LaTeXed}
solutions.


\begin{homeworkExercise}{Simple Samplers.}
  Suppose $X$ is a random variable which takes on values in the
  interval $[0, 1]$; let $\E[X] = c$. Initially, you don’t know
  anything about $c$, or about the probability distribution of
  $X$. However, you are given a black-box that every time you query
  it, it gives you an independent random sample drawn according to
  $X$. You want to estimate $c$.  A natural scheme is: sample from the
  black-box $N$ times—call these samples $X_1, X_2, \ldots, X_n$— and
  return the empirical mean
  $\hat{c} := \frac{1}{N}\sum_{i = 1}^N X_i$.  The natural question
  is: how big does the number of samples $N$ have to be so that
  \begin{align}\P[|\hat{c} - c| \leq \epsilon] \geq 1 - \delta \label{eqn:conf}\end{align}
  I.e., you
  want to be within error $\epsilon$ with confidence $1 - \delta$.  
  
  \begin{enumerate}[label=\alph*.]
    \item  Use
      Chebyshev's inequality to show that $N = O\left(\frac{ 1}{ \epsilon^2\delta} \right)$
  samples suffice to Ensure ~\ref{eqn:conf}.
  Hence, to get $\delta = \frac{1}{n^k}$ for some value $n$, we would
  take $O\left(\frac{n^k}{\epsilon^2}\right)$ samples.
  \item Hoeffding's bound says the following
    (you don’t have to prove this, of course):

    \begin{theorem}  Suppose $Y_1,
  Y_2,\ldots Y_T$ are independent $[0, M]$-bounded random variables, and
  define $Y := \sum^T_{ t=1} Y_t$ be their sum. Let $\mu = \E[Y ]$. Then
  \[\P[Y \geq \mu + \lambda] \leq \exp( -\frac{ \lambda^2}{  M(2\mu+\lambda)})\]
  \[\P[Y \leq \mu - \lambda] \leq \exp(\frac{ - \lambda^ 2}{ 2M \mu}).\]
  \end{theorem}
  Use
  Hoeffding's bound to show that $N = O\left(\frac{1}{\epsilon^2}\ln \frac{1}{\delta} \right)2$ samples are
  sufficient. Hence, to get $\delta = 1/n^k$ we would take $O\left( \frac{k \log n}{\epsilon^2}\right )$
  samples.
\end{enumerate}
\end{homeworkExercise}

\begin{homeworkExercise}{Matrix Multiplication is Useful.}
  Given an undirected simple graph $G = (V, E)$, a triangle is just a
  clique of size $3$; i.e., $3$ vertices such that all $3$ edges are
  present. Give algorithms for the following problems:
    \begin{enumerate}[label=\alph*.]
    \item Find a triangle in $G$ in time $n^\omega$
      \item Find a $3k$-clique in a graph in time $n^{k \omega}$.
      \item (Bonus) Find a triangle in $G$ in time $m^{1.5}$. (This
        one is slightly harder—not an exercise— and does not use
        matrix multiplication.)
    \end{enumerate}
    
    \end{homeworkExercise}

    \begin{homeworkExercise}      {Low-Diameter Decompositions for Simple Graphs.}
      Recall the concept of a $\beta$-low-diameter
decomposition, which given graph $G$ and distance $D$, randomly breaks the graph
into pieces of max-distance $D$, such that each pair $x, y$ is separated with probability at most
$\beta \cdot \frac{d_G(x,y)}{D}$
\begin{enumerate}[label=\alph*.]
\item Show that if each edge $(x, y) \in E(G)$ is cut with probability $\beta \cdot \frac{d_G(x,y)}{D}$, , then so is any pair
$x, y \in V$ . Hence, if the graph only has unit-weight edges, each edge can be cut with
probability at most $\beta/D$.
\item  Show that (i) any path graph has an LDD with $\beta = 1$, (ii) any tree with $\beta = 2$, and (iii)
the standard $k$-dimensional $\left( n^{\frac{1}{k}} \times \ldots \times n^{\frac{1}{k}}\right)$-grid  with $\beta = k$.
\end{enumerate}   
      \end{homeworkExercise}

      
      \begin{homeworkExercise}{Approximation via Randomized Simplification}
        In Lecture \#5 we encountered low-stretch spanning trees. Now
        we use this to solve the following $k$-median problem on
        general metrics: given an $n$-point metric $(X, d)$, and a
        number $k$, find a set $C$ of size $k$ that minimizes
        $\Phi_d(C) : = \sum_{v\in X}d(v, C)$.We outlined an algorithm
        for the case of distances $d_T$ on a tree $T$. We now show now
        to solve this problem (approximately) on a general $n$-point
        metric $(X, d)$.
\begin{enumerate}[label=\alph*.]
\item Use a low-stretch spanning tree
        constuction on metric $(X, d)$ to sample a tree $T$ from $D$.  Show
        that if $C_X$ is the optimal solution on metric $X$, and $C_T$ is the
        optimal solution on the tree $T$, then \[\E[\Phi_{d_T} (C_T )] \leq \E[\Phi_{d_T}
          (C_X)] \leq \alpha \times  \Phi_d(C_X).\]
        \item Observe that you can solve the
        clustering problem optimally on $T$ to find the centers $C_T$ using
        HW1. (Nothing to show here.)
      \item Show that $\Phi_d(C_T ) \leq \Phi_{d_T} (C_T )$.
        \item
        Deduce that the expected cost of the solution $C_T$ is at most
        $O(\log n)$ times $\Phi_d(C_X)$, the cost of the optimal
        solution. (Hence you have output an “$O(\log n)$-approximate
        solution” to this NP-hard problem.)
        \item Look at your analysis
        in the previous part and outline for what kinds of problems on
        metrics does this method apply. E.g., does it work for the TSP
        on a metric? How about the  $K$-center problem? Or the $K$-means
        problem which wants to minimize $\Psi_d(C) := \sum_{v\in X}(d(v, C))^2$ . Why
        or why not?
        \end{enumerate}
\end{homeworkExercise}



\begin{homeworkExercise}({Approximate MSTs and Shortest Paths} Show
  that the expected weight of a low-stretch spanning tree as above is
  at most $O(\alpha)$ times the MST.
\end{homeworkExercise}

\begin{homeworkProblem}{Another Halving Trick.}
  Suppose you have an algorithm $A$ that, given a directed graph
$G$ with integer edge lengths we satisfying $w_e \geq -1$, returns feasible potentials $\Phi : V \to \R$
in time $T$ (or shows there is a negative cycle). We will show that you can use this to find
feasible potentials for graphs with integer weights $w_e \geq -M$ in time $O((T + m) \log M)$ (or
shows there is a negative cycle).
\begin{enumerate}[label=\alph*.]
\item Define weights $w^{'}_e = \lceil w_e/2\rceil $. Show that if there are no negative cycles with respect to
  $w$, there are none w.r.t. $w'$.
  \item Suppose $\Phi': 
 V \to \R$ gives feasible potentials w.r.t. $w'$. Then show that the reduced cost

 $\hat{w}_{uv} := 2\Phi'_u + w_{uv} - 2\Phi'_v$ satisfies
 $\hat{w}_e \geq -1$ for all edges $e$.
 \item Now use $A$ to find
 feasible potentials w.r.t. $\hat{w}$. How does this give feasible potentials
 w.r.t. $w$? Show that if $w_e \geq -M$, we take $O((T + m)\log M)$ time.
 \end{enumerate}
\end{homeworkProblem}

\begin{homeworkProblem}{ Short on Average.}
  Given a directed graph $G = (V, A)$ with possibly negative
  edge-weights ${w_a}_{a\in A}$, define the weight-ratio of a directed
  cycle $C$ to be $\rho(C) = \frac{\sum_{ a\in C} w_a}{ |C|}$ .  We want to
  find a cycle $C$ with minimum weight-ratio. Denote this by
  $\rho^* (G) := \min_{\textrm{cycles }C} \rho(C)$.  (You may
  assume that edge-weights are integers in the range $[-M, M]$.)
    \begin{enumerate}[label=\alph*.]
      \item  (Nothing to submit.) Observe
    that you can check if a graph has a negative-weight cycle, using
    Bellman-Ford-Moore. (If all vertices are reachable from s this is
    immediate, else what will you do?)
  \item Show how to use Bellman-Ford-Moore to find a cycle of zero
    weight assuming there is no negative-weight cycle. (Hint: Recall
    how you used B-F to compute feasible potentials in Johnson’s
    algorithm from Lecture \#4.)
  \item
    Observe:$ \rho^* (G) = max\{\alpha \in Q | G\textrm{ with arc
      weights $(w_a - \alpha)$ has no negative cycle }\}$.  Use this
    observation to compute $\rho^* (G)$, and also find a cycle
    $C$ of this weight ratio. Your algorithm should run in time
    $O(mn(\log M + \log n))$ time.  Caveat: binary search over a range
    of $K$ integers takes $O(\log K)$ time. But binary search over
    rationals or reals may not terminate.
    \end{enumerate}
  \end{homeworkProblem}

  \begin{homeworkProblem}{Sparsification in Dynamic Algorithms.}
    In this problem we develop an approach to
transform an $O(m^{2/3}
)$ update time algorithm for dynamic connectivity into an $O(n^{
2/3}\log n)$
one. For all parts, assume all edge sets are over vertex set $V$ with $|V | = n$, and all graphs are
simple. Let $K = \lceil \frac{n - 1}{2}\rceil$.
\begin{enumerate}[label=\alph*]
\item For two edge sets $E, E'$ with spanning forests $F, F'$
respectively, show that any spanning
forest $F''$
 of $F \cup F'$
is also a spanning forest of $E \cup E'$.
\item A sparse partition of edge set $E$ is a partition into $K$ sets
  ($E^1 , E^2 , \ldots, E^K$), where some prefix $E^1 ,\ldots, E^k$
  all have exactly $n$ edges, $E^{k+1}$ is potentially partially full,
  and $E^{k+2}$ onwards are empty.  Suppose $E$ has sparse partition
  ($E^1, E^2, \ldots E^K)$. Suppose $E'$ is obtained from $E$ by
  either adding or removing one edge. Show how to get a sparse
  partition for $E'$ such that $(E')^i = E^i$ for all but (at most)
  two indices $i$, and for these two indices $i$ the size of the
  symmetric difference $|(E' )^i\Delta E^i | \leq 2$.
\item Define a complete binary tree with leaves $\{1, 2, . . . ,
  K\}$. Each internal node corresponds to some interval $i \ldots j$,
  define $E^{ij} := \cup_{s = i }^j E^s$ , and $F^{ij}$ be a spanning
  forest on $E^{ij}$ . This defines a collection of $O(K)$ spanning
  forests.  Given a sparse partition $(E^1 , E^2 , . . . , E^K)$ for
  $E$, and an associated collection of spanning forests, suppose we
  add or remove a single edge from one of the $E^i$ s. Show how to
  maintain the resulting collection of spanning forests with total
  update time $O(n^{ 2/3} \log n)$.  Infer that we can maintain
  dynamic connectivity in $O(n^{ 2/3}\log n)$ time.
  \end{enumerate}
\end{homeworkProblem}

    \begin{homeworkProblem}{Sparse Spanners.} Given a graph $G$ with edge lengths $\ell_e$, a subgraph $H$ is a spanner with
stretch $\gamma\geq 1$ if for every edge $(x, y) \in E(G),
d_H(x, y) \leq\gamma \cdot d_G(x, y)$.
\begin{enumerate}[label=\alph*.]
\item Use the triangle inequality to show that for all $x, y \in V$ , even if $(x, y)$ is not an edge,
$d_H(x, y) \leq\gamma \cdot d_G(x, y)$.
Clearly if $H = G$, we can set $\gamma = 1$. The goal is to find $H$ with few edges, such that $\gamma$ is also
small. Let’s give two different constructions of good spanners.
\item \emph{Approach \#1}. Consider the following randomized process: sample $t = 4 \log n$ trees
$T_1, T_2,\ldots, T_t$ from an $\alpha$-stretch probabilistic embedding into subtrees. Let $H$ be the
union of all these edges.
\begin{enumerate}
  \item
 Show that for any fixed edge $(x, y) \in E(G)$,
\[P[d_H(x, y)\geq 2\alpha d_G(x, y)] \leq 2^{-t}.\]

(Hint: for any single value of $i$, bound $\P[d_{T_i}
(x, y)\geq 2\alpha d_G(x, y)]$.)
\item Use the results quoted in Lecture to show that with probability $1 -\frac{1}{
n^2}$ , the graph
$H$ is an $O(\log n \log \log n)$-stretch spanner with $O(n \log n)$ edges.

\end{enumerate}
\item\emph{Large-girth graphs are Sparse.} On a seemingly different
  note, here is a useful graphtheoretic fact. Define the girth of a
  graph $G$ is the length of the shortest cycle in $G$. We will show
  that any graph $G$ with $m$ edges and $n$ nodes, and girth strictly
  more than $g$ must have $m \leq O(n + n^{ 1+\frac{1}{\lfloor g/2\rfloor}})$.
  \begin{enumerate}
    \item
  $G$ has average degree $\hat{d} := \frac{2m}{n}$ . Show that there exists a subset $S\subseteq V$ such that the induced subgraph $H := G[S]$ has minimum degree at
  least $\hat{d}/2$. [Hint: drop some low-degree vertices.]
  \item . For this
  graph $H$ and any vertex $v \in H$, show that the number of distinct
  vertices at distance at most $\lfloor g/2\rfloor$ from $v$ is at least $( \hat{d}/2 -
  1)^{\lfloor g/2\rfloor}$ .
  \item Show that the number of edges in the original graph
    $G$, i.e., $m \leq O(n + n^{ 1+\frac{1}{\lfloor g/2\rfloor}} )$.
  \end{enumerate}
  \item (d) \emph{Approach \#2}. Now consider
  a variant of Kruskal’s algorithm for $\alpha\geq 1$. Consider the
  edges of $G$ in increasing order of lengths $e_1, e_2, \ldots,
  e_m$. Initialize $H_0 = \emptyset$. When considering edge $e_i = (x, y) \in E(G)$,
  if the current distance $d_{H_{i-1}} (x, y) \leq\alpha d_G(x, y)$, then
  discard $e$ (i.e., set $H_i \gets H_{i-1}$, else set $H_i \gets H_{i-1} \cup \{e_i\}$.
  \begin{enumerate}
  \item (Do not submit.) Show that if we set $\alpha = n - 1$, then
    you will get Kruskal’s algorithm. Also, observe that by
    construction, the graph $H$ at the end of the process is an
    $(n - 1)$-stretch spanner. (In fact, an $(n - 1)$-stretch spanning
    tree.)
  \item  If we set $\alpha = O(\log n)$,
  use (c) with $g = O(\log n)$ to show the final graph $H$ is a $O(\log
  n)$-stretch spanner with $O(n)$ edges.
\end{enumerate}
\end{enumerate}
    \end{homeworkProblem}

\end{document}
